problem:  
Let \(V^n \subset S^{n+k}\) be a stationary integral \(n\)-varifold whose regular part is a smooth minimal submanifold, and assume that each tangent cone \(C_p = \mathbb{R}_+ \times L_p\) at a singular point \(p \in \mathrm{sing}\,V\) is multiplicity one, with \(L_p \subset S^{n+k-1}\) a smooth minimal submanifold (the *link*).  

For each link \(L_p\), let \(\mathcal{J}_{L_p}\) denote the **Jacobi operator** acting on sections of the normal bundle of \(L_p\) in \(S^{n+k-1}\):
\[
\mathcal{J}_{L_p} = \Delta^{\perp} + |A_{L_p}|^2 + \mathrm{Ric}_{S^{n+k-1}}(\nu,\nu),
\]
where \(\Delta^{\perp}\) is the connection Laplacian on normal sections and \(\nu\) is a local orthonormal frame of the normal bundle.  
Let \(\mathrm{ind}(L_p)\) denote the number of negative eigenvalues of \(\mathcal{J}_{L_p}\) (counted with multiplicity).

**Question:**  
Assume that 
\[
\mathrm{ind}(L_p) \le N
\]
for some fixed integer \(N\), uniformly for all tangent cone links \(L_p\) of \(V\).  

Does this *uniform index bound* on the links enforce a corresponding improvement in the regularity of \(V\)?  
In particular, does there exist a universal \(d_{n,k,N} < n\) such that
\[
\dim_{\mathcal{H}}(\mathrm{sing}\,V) \le d_{n,k,N}?
\]
If not, can one construct stationary minimal varifolds in the round sphere whose tangent cone links all have bounded index but whose singular sets accumulate with large Hausdorff dimension—potentially approaching the conjectural maximal dimension of \((n-7)\)?  

Equivalently, does bounded spectral instability of tangent cones impose any measure-theoretic or dimensional rigidity on singular sets in the positively curved ambient sphere?

Why is it a "good" problem:  
This refined problem formulates a sharp question connecting *spectral stability* (through the Morse index of tangent cone links) with *stratified regularity theory* (through Hausdorff dimension of singular sets)—an interplay central yet unexplored in current GMT.  

The use of index, rather than arbitrary spectral sums, gives a geometric and analytically natural invariant: it measures how far each cone link deviates from being stable. The conjectured link between a uniform index bound and upper bounds on singular dimension would represent a new kind of **index–dimension principle**, potentially extending the spirit of Federer–Almgren’s regularity theory into a spectral regime.  

Positive resolution could reveal a deep new mechanism by which ambient curvature and spectral constraints enforce geometric regularity, merging stability analysis, PDE spectral theory, and the microlocal structure of singularities.  
A negative answer, via explicit failures, would clarify the limits of spectral control in the presence of positive curvature.  

The statement is simple and fully geometric—expressed only via minimal links, their stability indices, and singular set dimension—but its resolution demands substantial new insights bridging curvature, spectral analysis, and varifold regularity. It thus exemplifies the “narrow entrance, profound depth” character of a truly good research problem.